\chapter{Traditional User Pairing Methods}
\label{chap:traditional}

Having established the system model and channel generation framework in the previous chapter, we now turn our attention to traditional baseline user pairing strategies for NOMA systems. These methods represent the current state-of-the-art in practical NOMA deployments and optimization research, spanning from simple heuristic approaches requiring minimal computation to sophisticated graph-theoretic algorithms that guarantee globally optimal solutions. Understanding these traditional methods is essential for establishing performance benchmarks against which our proposed deep learning approach will be evaluated, and for appreciating the trade-offs between computational complexity and solution quality that motivate machine learning research in this domain.

We present four representative pairing strategies that collectively cover the design space from greedy heuristics through optimal combinatorial optimization. Static pairing employs the simplest possible strategy of matching strongest users with weakest users to maximize channel disparity. Balanced pairing moderates this approach by dividing users into upper and lower halves by channel quality and pairing within reasonable bounds. Blossom maximum-weight matching formulates the problem as general graph optimization and applies Edmonds' algorithm to find globally optimal solutions. Finally, bipartite graph matching with proportional fairness weighting provides a middle ground that exploits the natural strong-weak structure of NOMA while incorporating fairness objectives. For each method, we provide the algorithmic framework, computational complexity analysis, implementation details, and discuss advantages and limitations in the context of practical NOMA deployment.

\section{Static Pairing Method}

The static pairing method represents the simplest heuristic approach to user clustering in NOMA systems. This method operates on a straightforward principle: pair the user with the strongest channel gain with the user having the weakest channel gain, continuing this process iteratively until all users are paired.

\subsection{Algorithm Description}

\textbf{Rationale:} Maximizes channel disparity $|h_j|^2 / |h_i|^2$, ensuring SIC feasibility and allocating high power to cell-edge users.

\begin{algorithm}[H]
\caption{Static Pairing}
\label{alg:static}
\begin{algorithmic}[1]
\Require Sorted users $1 \leq \cdots \leq N$ by channel gain
\Ensure Pairings $\mathcal{M}$

\State $\mathcal{M} \gets \emptyset$
\State $left \gets 1$, $right \gets N$

\While{$left < right$}
    \State $\mathcal{M} \gets \mathcal{M} \cup \{(left, right)\}$
    \State $left \gets left + 1$
    \State $right \gets right - 1$
\EndWhile

\State \Return $\mathcal{M}$
\end{algorithmic}
\end{algorithm}

\subsection{Complexity Analysis}

\begin{table}[H]
\centering
\begin{tabular}{ll}
\toprule
\textbf{Operation} & \textbf{Complexity} \\
\midrule
Sorting users & $O(N \log N)$ \\
Pairing iteration & $O(N)$ \\
\textbf{Total} & $\boldsymbol{O(N \log N)}$ \\
\bottomrule
\end{tabular}
\caption{Computational Complexity of Static Pairing}
\end{table}

\subsection{Advantages and Limitations}

The advantages of static pairing are compelling from an implementation perspective. The algorithm is extremely simple to implement, requiring minimal computational overhead suitable for resource-constrained base stations. By construction, it guarantees SIC feasibility through the largest possible channel gain disparity, and it naturally allocates higher power to cell-edge users who need it most.

However, these benefits come with significant limitations. The method completely ignores spatial correlation and angular positions of users, which can lead to pairing users with similar angular coordinates and consequently high inter-pair interference. For non-uniform user distributions, the fixed strongest-weakest pairing rule becomes suboptimal. Additionally, power allocation often wastes resources by giving excessive power to the stronger user in each pair who may not need it for achieving target rates.

\section{Balanced Pairing Method}

Balanced pairing addresses some limitations of static pairing by adopting a more moderate approach. Instead of pairing the extremes (strongest with weakest), this method divides users into upper and lower halves based on channel gain, then pairs corresponding ranks from each half.

\subsection{Algorithm Description}

\textbf{Rationale:} Improves fairness by avoiding extreme pairings while maintaining reasonable SIC feasibility.

\begin{algorithm}[H]
\caption{Balanced Pairing}
\label{alg:balanced}
\begin{algorithmic}[1]
\Require Sorted users $1 \leq \cdots \leq N$
\Ensure Pairings $\mathcal{M}$

\State $mid \gets \lfloor N/2 \rfloor$
\State $\mathcal{M} \gets \emptyset$

\For{$i = 1$ to $mid$}
    \State $\mathcal{M} \gets \mathcal{M} \cup \{(i, mid + i)\}$
\EndFor

\State \Return $\mathcal{M}$
\end{algorithmic}
\end{algorithm}

\subsection{Complexity Analysis}

The computational complexity remains the same as static pairing: $O(N \log N)$ dominated by the sorting operation.

\subsection{Advantages and Limitations}

Balanced pairing offers better fairness compared to static pairing through more uniform rate distribution across paired users. The channel gain differences within pairs are smaller but still sufficient for SIC in most scenarios. The method remains computationally efficient with the same O(N log N) complexity.

The primary limitation is that the fixed pairing rule doesn't adapt to actual channel statistics or spatial distribution of users. For clustered user distributions where users in the upper half may be spatially close, balanced pairing may violate SIC constraints more frequently than static pairing. The method also ignores angular separation, potentially creating high inter-pair interference in certain deployment geometries.

\section{Blossom Maximum-Weight Matching}

Blossom matching represents the gold standard for optimal user pairing in NOMA systems. Unlike the heuristic methods described previously, Blossom formulates pairing as a maximum-weight matching problem on a general graph and solves it to global optimality.

\subsection{Graph Formulation}

The user pairing problem is modeled as a general graph $G = (V, E)$ with the following structure:

\begin{itemize}
    \item \textbf{Vertices:} Users $V = \{1, \ldots, N\}$
    \item \textbf{Edges:} $(i,j) \in E$ if pairing is feasible (SIC + angular constraints)
    \item \textbf{Weights:} $w_{ij} = R_{sum}(i,j,\alpha^*)$ (optimized sum-rate)
\end{itemize}

\textbf{Objective:} Find maximum-weight matching:
\begin{equation}
\max \sum_{(i,j) \in \mathcal{M}} w_{ij}
\end{equation}
subject to each node appearing in at most one edge.

\subsection{Edmonds' Blossom Algorithm}

The Blossom algorithm, developed by Jack Edmonds in 1965, solves maximum-weight matching on general graphs through three key mechanisms:

\begin{enumerate}
    \item Finding augmenting paths in a weighted graph
    \item Handling odd cycles ("blossoms") through contraction
    \item Guaranteeing global optimality
\end{enumerate}

\subsection{Implementation}

\begin{lstlisting}[language=Python, caption={Blossom Matching Implementation}]
import networkx as nx

def blossom_matching(h_values, angles, P, sigma2, sic_th_db=8):
    """Perform Blossom maximum-weight matching."""
    N = len(h_values)
    G = nx.Graph()
    G.add_nodes_from(range(N))
    
    # Add weighted edges for feasible pairs
    for i in range(N):
        for j in range(i+1, N):
            # Check SIC feasibility
            h_ratio_db = 10 * np.log10(h_values[j] / h_values[i])
            if h_ratio_db < sic_th_db:
                continue
                
            # Check angular separation
            angle_diff = abs(angles[i] - angles[j])
            angle_diff = min(angle_diff, 2*np.pi - angle_diff)
            if np.degrees(angle_diff) < 25:
                continue
            
            # Compute optimal sum-rate
            alpha_opt = optimize_power(h_values[i], h_values[j],
                                      P, sigma2, sic_th_db)
            R_sum = compute_sum_rate(h_values[i], h_values[j],
                                    alpha_opt, P, sigma2)
            
            G.add_edge(i, j, weight=R_sum)
    
    # Find maximum-weight matching
    matching = nx.max_weight_matching(G, maxcardinality=False)
    return list(matching)
\end{lstlisting}

\subsection{Complexity Analysis}

\begin{table}[H]
\centering
\begin{tabular}{ll}
\toprule
\textbf{Operation} & \textbf{Complexity} \\
\midrule
Edge weight computation & $O(N^2)$ pairs × $O(\log(1/\epsilon))$ \\
Blossom algorithm & $O(N^3)$ worst-case \\
\textbf{Total} & $\boldsymbol{O(N^3)}$ \\
\bottomrule
\end{tabular}
\caption{Computational Complexity of Blossom Matching}
\end{table}

For $N=500$ users, runtime reaches approximately 44.7 seconds on standard hardware, making real-time deployment challenging.

\subsection{Optimality Guarantees}

Blossom provides \textbf{globally optimal} matching given fixed power allocations $\alpha^*$ for each pair and accurate sum-rate calculations. We use Blossom as the \textbf{performance upper bound} for comparison with all other methods. The optimality comes at the cost of O(N³) computational complexity, which becomes prohibitive for large N in real-time scenarios.

\section{Bipartite Graph Matching}

Bipartite graph matching represents a middle ground between heuristic methods and general graph optimization, exploiting the natural structure of NOMA user pairing to achieve near-optimal performance with improved computational efficiency.

\subsection{Strong-Weak Partitioning}

Divide users into two sets based on channel gain:
\begin{itemize}
    \item $V_1$: "Weak" users (lower 50\% by channel gain)
    \item $V_2$: "Strong" users (upper 50\%)
\end{itemize}

Create bipartite graph $G = (V_1 \cup V_2, E)$ where edge $(i,j)$ exists if $i \in V_1$, $j \in V_2$, and feasibility constraints hold.

\subsection{Proportional Fairness Weighting}

Instead of raw sum-rates, use proportional fairness weights:
\begin{equation}
w_{ij}^{PF} = \log(R_i + R_j + \epsilon)
\end{equation}

This logarithmic transformation balances throughput and fairness, preventing monopolization by high-rate pairs.

\subsection{Maximum-Weight Matching on Bipartite Graphs}

We use NetworkX's \texttt{max\_weight\_matching} algorithm which efficiently solves maximum-weight matching on bipartite graphs. While this has the same $O(N^3)$ complexity as the Hungarian algorithm, the bipartite structure provides practical speedups.

\subsection{Implementation}

\begin{lstlisting}[language=Python, caption={Bipartite PF Matching using NetworkX}]
import networkx as nx

def build_bipartite_pf_matching(h_values, angles, theta_min_deg=25):
    """
    Build bipartite graph (weak <-> strong) with PF weights and 
    angular guard, then run max-weight matching.
    """
    N = len(h_values)
    sorted_idx = np.argsort(h_values)
    theta_min_rad = np.deg2rad(theta_min_deg)
    
    # Split into weak and strong users
    weak = list(sorted_idx[:N//2])
    strong = list(sorted_idx[N//2:])
    
    G = nx.Graph()
    
    for i in weak:
        for j in strong:
            # Angular guard (shortest arc)
            angle_diff = abs(angles[i] - angles[j])
            angle_diff = min(angle_diff, 2*np.pi - angle_diff)
            if angle_diff < theta_min_rad:
                continue
            
            # SIC guard
            h1, h2 = h_values[i], h_values[j]  # i weak, j strong
            h_ratio_db = 10 * np.log10(h2 / h1)
            if h_ratio_db < 8:  # SIC threshold
                continue
            
            # Compute PF weight
            P1, P2, R1, R2, R_sum = calc_pair_rate(h1, h2)
            pf_weight = np.log(R1 + 1e-12) + np.log(R2 + 1e-12)
            
            if np.isfinite(pf_weight):
                G.add_edge(i, j, weight=pf_weight)
    
    # Find maximum-weight matching
    matching = nx.max_weight_matching(G, maxcardinality=True)
    return list(matching)
\end{lstlisting}

\subsection{Complexity and Advantages}

\textbf{Complexity:} Same as Blossom at $O(N^3)$, but with tighter practical constants due to bipartite structure.

\textbf{Advantages:}
\begin{itemize}
    \item Enforces strong-weak pairing (natural for NOMA)
    \item PF weighting balances throughput and fairness
    \item Slightly faster than Blossom in practice
    \item Improves Jain's fairness index by 10-15\%
\end{itemize}

\section{Power Allocation Subroutine}

All methods require solving the power allocation problem for each candidate pair to determine optimal NOMA power coefficients.

\subsection{Optimization Problem}

For a given pair $(i,j)$, find:
\begin{equation}
\alpha^* = \argmax_{\alpha \in [\alpha_{min}, 1)} R_{sum}(i,j,\alpha)
\end{equation}
subject to SIC constraint.

\subsection{Bisection Search Algorithm}

\begin{algorithm}[H]
\caption{Power Allocation via Bisection}
\label{alg:power}
\begin{algorithmic}[1]
\Require Channel gains $h_i, h_j$, power $P$, noise $\sigma^2$, SIC threshold $\gamma_{th}$
\Ensure Optimal power split $\alpha^*$

\State Compute $\alpha_{min}$ satisfying SIC constraint
\State $a \gets \alpha_{min} + \epsilon$, $b \gets 1 - \epsilon$
\State $tol \gets 10^{-4}$, $max\_iter \gets 80$

\For{$iter = 1$ to $max\_iter$}
    \State $\alpha_{mid} \gets (a + b) / 2$
    \State $R_{sum}(mid) \gets $ Evaluate sum-rate at $\alpha_{mid}$
    \State $R_{sum}(left) \gets $ Evaluate at $\alpha_{mid} - tol$
    \State $R_{sum}(right) \gets $ Evaluate at $\alpha_{mid} + tol$
    
    \If{$R_{sum}(right) > R_{sum}(mid)$}
        \State $a \gets \alpha_{mid}$
    \ElsIf{$R_{sum}(left) > R_{sum}(mid)$}
        \State $b \gets \alpha_{mid}$
    \Else
        \State \textbf{break} \Comment{Converged}
    \EndIf
\EndFor

\State \Return $\alpha_{mid}$
\end{algorithmic}
\end{algorithm}

\subsection{Computational Cost}

Per pair: $O(\log(1/tol)) \approx 20$ iterations × 3 function evaluations = 60 operations.

For $N$ users: $O(N^2)$ candidate pairs, total $O(N^2 \log(1/tol))$ for power optimization across all methods.

\section{Comparative Summary}

Table \ref{tab:trad_methods} provides a comprehensive comparison of all traditional methods across key performance dimensions.

\begin{table}[H]
\centering
\caption{Comparison of Traditional Methods}
\label{tab:trad_methods}
\begin{tabular}{lcccc}
\toprule
\textbf{Method} & \textbf{Complexity} & \textbf{Optimality} & \textbf{Fairness} & \textbf{Constraints} \\
\midrule
Static & $O(N \log N)$ & Low & Low & Implicit SIC \\
Balanced & $O(N \log N)$ & Medium & Medium & Implicit SIC \\
Blossom & $O(N^3)$ & \textbf{Optimal} & High & Explicit \\
Bipartite PF & $O(N^3)$ & Near-Optimal & \textbf{Highest} & Explicit \\
\bottomrule
\end{tabular}
\end{table}

\textbf{Key Observations:}
\begin{itemize}
    \item Heuristics (Static, Balanced) sacrifice 15-25\% throughput for speed
    \item Optimization methods (Blossom, Bipartite) achieve 92-100\% optimal but with $O(N^3)$ cost
    \item PF weighting improves fairness without significant throughput loss
    \item All methods struggle with $N > 50$ for real-time constraints (<1 ms)
\end{itemize}

\section{Computational Complexity Analysis and Comparison}

To provide concrete understanding of the computational burden imposed by each algorithm, we analyze both theoretical complexity and empirical runtime measurements. Table~\ref{tab:complexity_comparison} summarizes the key complexity characteristics.

\begin{table}[h]
\centering
\caption{Computational Complexity Comparison of Traditional Pairing Methods}
\label{tab:complexity_comparison}
\begin{tabular}{lccc}
\toprule
\textbf{Method} & \textbf{Time Complexity} & \textbf{Space Complexity} & \textbf{Optimality} \\
\midrule
Static Pairing & $O(N \log N)$ & $O(N)$ & Heuristic \\
Balanced Pairing & $O(N \log N)$ & $O(N)$ & Heuristic \\
Blossom Matching & $O(N^3)$ & $O(N^2)$ & Globally Optimal \\
Bipartite PF Matching & $O(N^3)$ & $O(N^2)$ & Optimal (bipartite) \\
\bottomrule
\end{tabular}
\end{table}

Static and balanced pairing achieve linearithmic $O(N \log N)$ complexity dominated by the sorting operation, making them suitable for real-time operation even with thousands of users. Their space complexity is linear $O(N)$, requiring storage only for user channel gains and pairing results.

Both graph-based methods exhibit cubic $O(N^3)$ time complexity, though this arises from different sources. For Blossom, the cubic term comes from augmenting path iterations in the matching algorithm itself. For bipartite matching, it comes from the Hungarian algorithm's dual update iterations. The space complexity is quadratic $O(N^2)$ due to explicit or implicit storage of the edge weight matrix.

The practical implications of these complexity differences are substantial. For $N = 12$ users, all methods execute in milliseconds, making complexity differences negligible. For $N = 100$ users, heuristics still execute in microseconds while graph methods require hundreds of milliseconds—noticeable but often acceptable. For $N = 500$ users representing dense deployments, heuristics execute in under 1 millisecond while Blossom requires 40-50 seconds, a difference of 5 orders of magnitude that clearly prohibits real-time Blossom operation.

This analysis reveals the fundamental motivation for machine learning approaches: if we can train a model offline using Blossom-generated optimal solutions as supervision, then deploy the trained model for real-time inference with complexity comparable to heuristics, we could achieve near-optimal performance with practical computational requirements. This is precisely the goal of our proposed Graph Neural Network framework presented in the following chapter.

\section{Summary and Insights}

This chapter has presented four traditional user pairing strategies representing the current state-of-the-art in NOMA research and practice. Static pairing offers maximum simplicity and efficiency but sacrifices optimization. Balanced pairing improves fairness while retaining efficiency. Blossom matching achieves global optimality at the cost of cubic complexity. Bipartite PF matching balances throughput, fairness, and computational cost through structured problem formulation.

Several key insights emerge from this analysis. First, a fundamental trade-off exists between solution quality and computational complexity—achieving theoretical optimality requires accepting impractical computation times for large-scale systems. Second, incorporating domain structure (the weak-strong nature of NOMA) into the problem formulation yields substantial benefits, as evidenced by bipartite matching's competitive performance. Third, fairness and efficiency are competing objectives that require explicit multi-objective optimization rather than hoping heuristics naturally balance them.

These traditional methods provide essential baselines for evaluating our proposed deep learning approach. Static and balanced pairing represent the lower complexity bound and heuristic performance levels we must exceed to justify machine learning's added sophistication. Blossom matching represents the upper performance bound we aspire to approach. Bipartite PF matching represents the current best practice balancing multiple objectives. In the next chapter, we present our Graph Neural Network framework designed to achieve Blossom-level performance with heuristic-level complexity, thereby bridging this critical gap for practical NOMA deployment.
